\chapter{一个变量：反馈怎样带来恢复或记忆？}
现在让产生率依赖蛋白质自身。单个动态变量只能沿浓度轴运动，但产生曲线的形状已经能决定系统是回到唯一表达水平，还是在不同长期状态之间作出选择。

\begin{figure}[H]\centering
\begin{tikzpicture}
\node[gene] (none) at (0,0) {$X$};\node at (0,-.8) {无反馈};
\node[gene] (pos) at (4,0) {$X$};\node at (4,-.8) {自激活};
\node[gene] (neg) at (8,0) {$X$};\node at (8,-.8) {自抑制};
\draw[act] (pos) to[loop above,min distance=12mm] (pos);
\draw[rep] (neg) to[loop above,min distance=12mm] (neg);
\end{tikzpicture}
\caption{单节点的三种符号拓扑。所有模型仍包含产生与降解；自边表示自身浓度对产生率的作用。}
\end{figure}
\section{自抑制加强恢复，为什么不会产生多个稳态？}
无反馈模型的产生率恒为 $\alpha$。自抑制把它改成随 $x$ 非增的 $F(x)$，例如 $F(x)=a/(1+(x/K)^n)$。浓度上升时，产生减少而降解增加，两种作用都将浓度拉回。

把这种直觉写成方程，令 $f(x)=F(x)-\gamma x$，假设产生率 $F(x)$ 非负、有界、非增且为 $C^1$。这里“有界”就是有一个有限上限，“非增”表示浓度增加时产生率不会上升，$C^1$ 则保证可以连续地计算斜率。$f$ 就严格递减，且 $f(0)\geq0$、$f(x)\to-\infty$。因此产生曲线与降解直线恰好相交一次，稳态 $x^*$ 唯一，允许它位于边界 $x^*=0$。

不仅稳态唯一，所有非负初值也都会趋于它。令浓度偏差 $e=x-x^*$，则
\[
\frac12\frac{\dd}{\dd t}e^2=e\{F(x)-F(x^*)\}-\gamma e^2\leq-\gamma e^2,
\]
因为 $F$ 非增，浓度偏差与产生率偏差的乘积不可能为正。积分可得
\[
|x(t)-x^*|\leq e^{-\gamma t}|x(0)-x^*|.
\]
式~\eqref{eq:bound}说明浓度不会无限增大，所以这个恢复速度的估计可以一直使用。若 $F'(x^*)<0$，小扰动满足 $\dot e\simeq(F'(x^*)-\gamma)e$，小偏差减小到原来约 $37\%$ 所需的时间为 $1/(\gamma-F'(x^*))$，短于无反馈的 $1/\gamma$。这把“自抑制促进恢复”落实为具体速率；它描述的是稳态附近的小偏差；输入变化很大时，还要直接计算整条时间轨迹。

\section{自激活让产生与降解出现三次交点}
自激活的方向相反：浓度上升也会提高产生率，偏差可能被放大。于是一个自然的问题是，这种放大能否使系统保留不同表达状态？取
\begin{equation}\dot x=a+b\frac{x^n}{K^n+x^n}-\gamma x.\label{eq:auto}\end{equation}
稳态是产生曲线与线性降解直线的交点。局部稳定条件为
\begin{equation}F'(x^*)-\gamma<0.\end{equation}
图~\ref{fig:auto} 的参数 $a=0.05,b=3.5,K=\gamma=1,n=4$ 产生三个稳态，外侧稳定，中间不稳定。关键不是正反馈必然双稳态，而是正反馈与足够的非线性产生多个交点。改变参数时，两个交点可能靠近并合在一起，再一起消失。合并的临界位置称为“折叠点”：此时产生曲线与降解直线不仅相交，斜率也相同，即 $F(x^*)=\gamma x^*$ 和 $F'(x^*)=\gamma$。
\begin{figure}[H]\centering\includegraphics[width=.74\linewidth]{figures/self_activation.pdf}\caption{自激活模型的产生率与降解率，实心点为稳定的稳态，空心点为不稳定的稳态。}\label{fig:auto}\end{figure}

先判断交点两侧的运动方向，再问交点何时会成对出现。对只有一个浓度变量的 $C^1$ 方程，若 $f'(x^*)<0$，则足够靠近稳态时 $f$ 在左侧为正、右侧为负，轨道向稳态运动；若 $f'(x^*)>0$ 则方向相反。如果交点处的斜率为零，就需要再看它左右两侧的实际变化方向，不能只凭这个斜率判断。对式~\eqref{eq:auto}，
\[
F'(x)=\frac{bnK^nx^{n-1}}{(K^n+x^n)^2}.
\]
当 $n>1$ 时该导数的唯一最大值出现在 $x_m=K((n-1)/(n+1))^{1/n}$：令 $h=x/K$，对 $h^{n-1}/(1+h^n)^2$ 求导，驻点条件为 $(n-1)-(n+1)h^n=0$。因此
\[
\max F'=\frac bK\frac{(n+1)^2}{4n}
\left(\frac{n-1}{n+1}\right)^{(n-1)/n}.
\]
若该值小于 $\gamma$，$f$ 严格递减，稳态唯一；超过 $\gamma$ 只是出现多交点的必要条件，仍需选择合适的 $a$。
令 $q=(x^*/K)^n$，联立折叠条件可参数化为
\[
b_{\rm fold}=\frac{\gamma K}{n}q^{(1-n)/n}(1+q)^2,\qquad
a_{\rm fold}=\frac{\gamma K}{n}q^{1/n}(n-1-q).
\]
这是由 $F'=\gamma$ 先解出 $b$，再代入 $a=\gamma x-bq/(1+q)$ 得到。对 $a>0$ 需 $0<q<n-1$。通常的折叠分岔还要求 $f''\ne0$，称为一个\emph{非退化条件}。“分岔”指参数变化使稳态或轨迹的类型发生改变；“非退化”在这里指相切处仍有明确的弯曲，两个交点的合并不是更特殊的平坦情形。若连弯曲程度也为零，交点怎样变化就需要进一步分析。

把条件落实到图中的参数，也能解析地确认这三个交点。取 $f(x)=0.05+3.5x^4/(1+x^4)-x$，有 $f(0)>0$、$f(0.1)<0$、$f(1)>0$、$f(4)<0$，所以三个互不相交的区间 $(0,0.1)$、$(0.1,1)$、$(1,4)$ 各含至少一个根。上面的计算说明 $f'$ 最多有两个零点，也就是说曲线最多改变两次升降方向。若 $f$ 有四个交点，相邻交点之间就至少需要三个斜率为零的位置，超过了这个上限。这用到了\emph{Rolle 定理}：一条可求导的曲线在两端取相同值时，中间至少有一处斜率为零。因此三个交点就是全部结果。两个极值分别严格低于、严格高于零，因此三个根均不在极值点处，斜率依次为负、正、负。中间的不稳定交点构成一道分界：初始浓度在它以下时趋于低表达态，在它以上时趋于高表达态。短暂扰动跨过此阈值可造成持久状态切换。

\section{有记忆，为什么仍不能自行振荡？}
两个稳定表达态已经足以构成记忆，却还不足以构成时钟。对\emph{一维自治方程} $\dot x=f(x)$，“一维”表示只有一个动态变量，“自治”表示变化率由当前状态决定、不显式依赖时间。如果外部条件固定、也没有延迟，那么同一个浓度 $x$ 总对应同一个变化方向。浓度在一个区间里一直上升，就不能在没有其他变化的情况下突然掉头下降。

可以把这一点完整地论证出来。假设有一条重复往返的轨迹，周期为 $T$。它在一个周期内一定有最高浓度 $x_m$。最高点处的变化率为零，所以 $f(x_m)=0$，也就是说这里已经是稳态。从这个浓度出发，始终停在 $x_m$ 就满足原方程；同一初始浓度又只能产生一条轨迹，因此原来的轨迹也只能停在那里，与“重复往返”矛盾。

所以，一个变量可以提供不同表达状态之间的选择，但在这些假设下不能自行维持振荡。要寻找持续振荡，需要增加动态环节，或者另行考虑延迟、随时间变化的输入等因素。
